\chaptertitle{第五章\quad 齐次坐标与仿射变换}

\sectiontitle{5.1\quad 平移——一个无法用$2\times2$矩阵表示的变换}

前四章我们学习的所有变换都能写成$x'=ax+by$，$y'=cx+dy$。但有一个最常见的操作——平移$(x,y)\to(x+a,y+b)$——上面这种形式写不出来。

为什么呢？因为平移在$x'$和$y'$的表达式里各多了一个常数项$a$和$b$。标准形式$x'=ax+by$要求$x'$是$x$和$y$的线性组合，不能有游离的常数。这暴露出一个基本限制：平移改变了原点，而线性变换（$ax+by$型）把原点固定在原处。平移后原点$(0,0)$去了$(a,b)$，不再是原点了。

为了把平移也统一进矩阵的框架，需要一个巧妙的办法。

\begin{exercise}
\item 验证：以下变换中哪些保持原点$(0,0)$不动？(a) 旋转$90^\circ$；(b) 平移$(2,-1)$；(c) 关于$x$轴反射；(d) 向右平移$3$再旋转$180^\circ$。
\item 为什么"保持原点不动"是能用$2\times2$矩阵表示的变换的必要条件？
\end{exercise}

\sectiontitle{5.2\quad 齐次坐标——给平面加一个维度}

我们不把平面上的点记作$(x,y)$，而是记作$(x,y,1)$。多出来的第三个坐标$"1"$只是一个占位符，不表示任何真实的空间方向，唯一的作用是让平移可以被"装"进矩阵乘法。

在齐次坐标下，向右平移$3$、向上平移$2$的操作写成：
\[
\begin{pmatrix}1&0&3\\0&1&2\\0&0&1\end{pmatrix}
\begin{pmatrix}x\\y\\1\end{pmatrix}
= \begin{pmatrix}x+3\\y+2\\1\end{pmatrix}
\]
你看，第一行把$x$坐标加了$3$，第二行加了$2$，第三行保持不变。平移就这样被$3\times3$矩阵表达出来了。

所有的$2\times2$线性变换（旋转、反射、拉伸、切变）在齐次坐标下也扩充为$3\times3$矩阵——只需把$2\times2$部分放在左上角，第三行第三列写$1$，其余位置写$0$：
\[
\begin{pmatrix}a&b&0\\c&d&0\\0&0&1\end{pmatrix}
\]

\begin{example}
将先沿$x$方向拉伸$2$倍，再向右平移$3$，用齐次坐标一步算出$(1,2)$的像。
\end{example}
\begin{solution}
线性部分$A=\begin{pmatrix}2&0\\0&1\end{pmatrix}$，平移$(3,0)$。
齐次矩阵：$\begin{pmatrix}2&0&3\\0&1&0\\0&0&1\end{pmatrix}$。
作用于$(1,2)$：$\begin{pmatrix}2\cdot1+0\cdot2+3\cdot1\\0\cdot1+1\cdot2+0\cdot1\\1\end{pmatrix}=\begin{pmatrix}5\\2\\1\end{pmatrix}$，即$(5,2)$。
\end{solution}

\begin{example}
用齐次坐标计算：先关于$x$轴反射，再平移$(-1,4)$，求点$(2,3)$的像。
\end{example}
\begin{solution}
反射矩阵$\begin{pmatrix}1&0\\0&-1\end{pmatrix}$，平移$(-1,4)$。齐次矩阵为：
\[
\begin{pmatrix}1&0&-1\\0&-1&4\\0&0&1\end{pmatrix}
\]
作用于$(2,3)$：$\begin{pmatrix}1\cdot2+0\cdot3+(-1)\cdot1\\0\cdot2+(-1)\cdot3+4\cdot1\\1\end{pmatrix}=\begin{pmatrix}1\\1\\1\end{pmatrix}$，即$(1,1)$。
\end{solution}

\begin{exercise}
\item 用齐次坐标计算：先旋转$90^\circ$再平移$(2,0)$，求点$(1,3)$的像。
\item 写出纯平移$(a,b)$的$3\times3$齐次坐标矩阵的一般形式。
\item 若一个$3\times3$齐次坐标矩阵的第三行不是$(0,0,1)$，它还能表示平面上的仿射变换吗？试分析。
\end{exercise}

\sectiontitle{5.3\quad 仿射变换的一般形式}

在齐次坐标下，一个任意仿射变换的矩阵长这样：
\[
\begin{pmatrix}
a&b&t_x\\
c&d&t_y\\
0&0&1
\end{pmatrix}
\]
左上角的$2\times2$子矩阵$A$是线性部分（旋转、反射等），右上角的$(t_x,t_y)$是平移部分。仿射变换就是"先线性变换，再平移"的总和。

按线性部分$A$的性质，可以给仿射变换分类。下面逐一用具体矩阵验证每类的特点。

\textbf{第一类：刚体运动。}$A$是正交矩阵，即$A^{\text T}A=I$，例如旋转矩阵或反射矩阵。此时$|\det A|=1$。

举个具体例子：$A=\begin{pmatrix}0&-1\\1&0\end{pmatrix}$（旋转$90^\circ$），$\vec t=(3,2)$。对应的$3\times3$齐次矩阵为$\begin{pmatrix}0&-1&3\\1&0&2\\0&0&1\end{pmatrix}$。这个变换先绕原点旋转$90^\circ$再向右上平移——相当于把一个图形在平面上刚性搬动。因为旋转不改变两点间的距离（$|A\vec v|=|\vec v|$对所有$\vec v$成立），平移也不改变距离，所以整个变换保持所有点间距离不变，面积绝对值也不变（$|\det A|=1$）。

\textbf{第二类：相似变换。}$A=kR$，其中$R$是旋转矩阵，$k>0$是缩放因子。例如$A=\begin{pmatrix}2&0\\0&2\end{pmatrix}$（整体放大$2$倍），$\vec t=(1,0)$。这个变换保持形状——原来的正方形放大后还是正方形，原来的圆放大后还是圆——只是大小变了。面积放大倍率为$|\det A|=k^2=4$。

\textbf{第三类：一般仿射变换。}$A$是任意可逆矩阵，例如$A=\begin{pmatrix}2&1\\0&3\end{pmatrix}$（含切变和不等比例伸缩）。$\vec t$任意。这类变换不保持角度和长度，但保持一个重要的几何性质：平行线在变换后仍然是平行线（你可以验证：两条原本平行的直线在$A$的作用下仍然平行，因为线性变换保持向量的线性关系）。

\textbf{第四类：降维变换。}$\det A=0$。例如$A=\begin{pmatrix}1&0\\0&0\end{pmatrix}$（向$x$轴投影），$\vec t$任意。整个平面被压扁到一条直线上，再平移——结果仍然是一条直线。这类变换不可逆，面积归零。

\sectiontitle{5.4\quad 变换的分解与复合}

任何一个仿射变换都可以唯一地分解为：先做线性部分$A$，再做平移$\vec t$。也可以反过来先平移再线性变换，这两种分解给出不同的矩阵表示，但都是同一个变换。

两个仿射变换的复合仍然是一个仿射变换。在齐次坐标中，复合就是矩阵乘法，非常方便。

\begin{example}
求两个仿射变换的复合：先做$T_1$（向右平移$2$），再做$T_2$（旋转$90^\circ$）。写出复合变换的齐次坐标矩阵，并计算点$(1,0)$的像。
\end{example}
\begin{solution}
$T_1$的矩阵：$\begin{pmatrix}1&0&2\\0&1&0\\0&0&1\end{pmatrix}$，$T_2$的矩阵：$\begin{pmatrix}0&-1&0\\1&0&0\\0&0&1\end{pmatrix}$。

复合（先$T_1$后$T_2$，$T_2$在前）：$T_2T_1=\begin{pmatrix}0&-1&0\\1&0&0\\0&0&1\end{pmatrix}\begin{pmatrix}1&0&2\\0&1&0\\0&0&1\end{pmatrix}=\begin{pmatrix}0&-1&0\\1&0&2\\0&0&1\end{pmatrix}$。

作用于$(1,0)$：$\begin{pmatrix}0&-1&0\\1&0&2\\0&0&1\end{pmatrix}\begin{pmatrix}1\\0\\1\end{pmatrix}=\begin{pmatrix}0\\3\\1\end{pmatrix}$，即$(0,3)$。

验证：先平移$(1,0)\to(3,0)$，再旋转$90^\circ$得$(0,3)$，一致。
\end{solution}

\begin{exercise}
\item 直接用齐次坐标矩阵乘法验证：先旋转$90^\circ$再向右平移$2$，与先平移$2$再旋转$90^\circ$，复合结果不同。指出两个复合变换分别把原点$(0,0)$变到了哪里——为什么不同？
\end{exercise}

\sectiontitle{本章小结}

仿射变换统一了线性变换和平移。核心技巧是齐次坐标——给平面上的点多加一维$(x,y,1)$，从而把平移"装进"$3\times3$矩阵的右上角。这使得旋转+平移（刚体运动）、伸缩+平移（相似变换）以及任意线性变换+平移都可以统一用$3\times3$矩阵表达和复合。按线性部分$A$的行列式，可以区分保持面积（$|\det A|=1$）、缩放面积（$|\det A|\neq1$）和降维（$\det A=0$）三类仿射变换。这套语言在计算机图形学中用来表达所有二维图形的移动、旋转和缩放。

\sectiontitle{课后练习}

\subsection*{A组\quad 基础巩固}
\begin{exercise}
\item 用齐次坐标矩阵计算：点$(2,3)$平移$(1,-2)$后的坐标。
\item 写出以下变换的$3\times3$齐次坐标矩阵：\\
(a) 先绕原点旋转$45^\circ$再向右平移$5$个单位；\\
(b) 先向$x$轴投影再向下平移$3$个单位。
\item 一个仿射变换的$3\times3$矩阵为$\begin{pmatrix}3&0&2\\0&2&1\\0&0&1\end{pmatrix}$。这个变换对点$(1,1)$作用后坐标是多少？它把面积放大了几倍？（提示：面积倍率只看线性部分$A$的$|\det A|$。）
\item 判断以下仿射变换是否保持面积：旋转+平移；投影+平移；切变+平移；反射+平移。
\item 写出先关于$x$轴对称再平移$(0,-2)$的齐次坐标矩阵，并计算$(3,4)$的像。
\end{exercise}

\subsection*{B组\quad 挑战提升}
\begin{exercise}
\item 一个仿射变换把$(0,0)$变成$(1,2)$，把$(1,0)$变成$(4,2)$，把$(0,1)$变成$(1,5)$。求出它的$3\times3$齐次坐标矩阵。
\item 二维平移矩阵是$3\times3$的，那么三维平移矩阵是几阶的？它长什么样？
\item 证明：两个刚体运动的复合仍然是刚体运动。
\item 一个仿射变换把$(0,0)$映射到$(2,-1)$，把$(1,0)$映射到$(5,1)$，把$(0,1)$映射到$(0,3)$。求该变换的$3\times3$齐次矩阵，并将它分解为"线性部分+平移部分"。
\item 用齐次坐标验证：先绕原点旋转$60^\circ$再向右平移$4$个单位的变换，与先向右平移$4$个单位再绕原点旋转$60^\circ$的变换，结果不同。写出两者的$3\times3$矩阵，并分别计算它们将$(1,1)$映射到了哪里（保留根号）。
\item 在计算机图形学中，一个矩形窗口的四个角点坐标为$(0,0),(3,0),(3,2),(0,2)$。若先将其绕自身中心旋转$180^\circ$再整体向右平移$5$个单位，写出对应的$3\times3$齐次矩阵并求变换后的四个角点坐标。
\end{exercise}

\newpage
